\documentclass[a4paper]{article}
\usepackage{luatexja}
\usepackage{amsmath}
\usepackage{amssymb}
\usepackage{geometry}
\usepackage{hyperref}

\geometry{margin=25mm}

\title{近接重力の斥力層としての強い力：\\
陽子のカオス、魔法数、遅い中性子、低温核融合}

\author{https://github.com/hypernumbernet}
\date{\today}

\begin{document}
\maketitle

\begin{abstract}
近接重力とは、二重時空理論の重力チャンネルが、その特性長 \(\ell\) に匹敵する間隔で働くときの力である。等スケール軌跡の上では、力はニュートンの力を干渉因子 \(\gamma_s\) で割ったものである。\(\gamma_s\) の符号が斥力である。質量は \(\gamma_s>0\) で引き合い、\(\gamma_s<0\) で斥け合う。最初の負の区間が斥力反転層である。その外側の縁は第一の節 \(x_1\) であり、\(\tan x=\coth x\) の \((\pi/4,\pi/2)\) における唯一の根で、数値は \(x_1=0.93755\)、\(r_1/\ell=0.53330\) である。内側の縁は第二の節で、\(r_2/\ell=0.12731\) にある。第一の節の両側は間隔を節へ向けて駆動し、等スケール軌跡の上では、その駆動は間隔を層の奥の引力の室までは運ばない。長さは荷電パイオンのコンプトン波長 \(\lambda_\pi\approx 1.41\,\mathrm{fm}\) である。実験室の符号は、それを \(56/75\,\mathrm{fm}\) と \(19/10\,\mathrm{fm}\) のあいだに厳密に置く。\(\tfrac25\,\mathrm{fm}\) での斥力は \(\ell>56/75\,\mathrm{fm}\) を強制し、\(1\,\mathrm{fm}\) での引力は \(\ell<19/10\,\mathrm{fm}\) を強制する。その長さのもとで、第一の節は \(14/19\) と \(213/280\,\mathrm{fm}\) のあいだに厳密にあり、ラピディティの壁は \(8/9\) と \(142/157\,\mathrm{fm}\) のあいだに厳密にある。

相対ラピディティ \(\varphi\) は \(\tanh\varphi=|v|/c\) で定まり、双対角を \(\ell/r\) に保ったまま、通常の近接ラピディティに加えられる。干渉因子は \(\cosh(\varphi/2)\,\gamma_s-\sinh(\varphi/2)\,\sigma\) になる。窓 \(r_1<r<2\ell/\pi\) の内部で、因子は閉じた閾値
\[
\varphi_*(x)=\ln\frac{\sin x+\cos x}{\sin x-\cos x}-2x,
\qquad x=\ell/(2r)
\]
において符号を変える。閾値より下ではチャンネルは引力であり、ラピディティが上がるほど引力は強まる。閾値より上ではチャンネルは斥力であり、外向きの力は閾値の直後で任意に大きい。半径 \(2\ell/\pi\) の外側では、有限なラピディティは斥力を生まない。したがって反転の外半径は \(\varphi\) とともに増し、静的な節の \(1.1937\) 倍である \(2\ell/\pi\) より厳密に小さいままである。小さなラピディティでは、節は位相において \(\varphi/(4\cosh^2 x_1)\) だけ動き、その係数は \(0.11528\) である。遅い中性子は、静的な節の上のその厚さの表皮まで引力の側に留まる。非負で、閾値より下にあるあいだ増え続けるラピディティは、\(x_1\) より厳密に小さい一つの位相で閾値に出会う。チャンネルが引力であるあいだ、自由な落下はその種の経路であるから、静的な節より前に閾値に出会う。チャンネルがなお引力であるあいだに節へ達するには、ラピディティが下がらなければならない。運動エネルギー \(1\,\mathrm{MeV}\)、\(10\,\mathrm{MeV}\)、\(100\,\mathrm{MeV}\) の中性子は、閾値に \(1.0056\,r_1\)、\(1.0175\,r_1\)、\(1.0517\,r_1\) で出会う。その閾値は入口の条件であり、捕獲と、中性子誘起核分裂と、核融合とに共通である。閾値より上では核のチャンネルはすでに斥力であり、三つとも始められない。これは分裂そのものではない。複合核の内部でラピディティが処分される仕方は導出されず、断面積が実験室で \(1/v\) に従って落ちる規則も導出されない。

陽子は二つで、電磁チャンネルも担う。核の前因子を \(\kappa\) と書くと、間隔を増す力は \(-\kappa/(\gamma\,r^2)+k e^2/(\gamma_E r^2)\) である。\(\gamma_s\) が零を下向きに横切る重力の節は間隔を引き寄せ、同じ種類の電磁の節は間隔を斥ける。一つのチャンネルの隣り合う節は、安定性が逆である。二つの長さが異なれば節は互い違いになり、引力と斥力は隣り合う室を占める。核の中では、異なる対が一つの瞬間に異なる室を占める。陽子に働く力は室の壁の上で非有界であり、一つの対を節の上に置く任意に小さい変位は、その力を任意に大きく変える。その敏感さが、カオス的な陽子力である。同定 \(\kappa=G m_p^2\) は、\(|\gamma|\) が因子 \(1.236\times 10^{36}\) を補わないかぎり、クーロン項に負ける。MeV の深さは入力である。

荷電の対に対しては、電磁の因子が正のままであるところでは、同じ閾値がクーロン項との争いを決める。閾値より上では核の項は外向きであり、正味の斥力はクーロン斥力を上回る。閾値より下では核の項は内向きであり、閾値で終わるラピディティの区間の上で、クーロンの前因子と \(\kappa\) とのあらゆる有限な比を上回る。より低いラピディティは、その始まりを静的な節の近くへ動かす。運動エネルギーが静止エネルギーの高々 \(10^{-10}\) であるとき、半径の分数的な広がりは \(5\times 10^{-6}\) より小さく、室温の核子または重陽子は、この種のものである。\(2\ell/\pi\) の外側では、大きなラピディティが核の引力を零へ追い、クーロン斥力を残す。自由な落下はなお節より前に閾値に出会うので、遅い荷電の入口とは、ラピディティが抑えられたままの経路である。その保持を一軸に沿って引き受ける相手は、正イオンの正八面体において、第一の電荷の節の外側にある戻す席である~\cite{dst-cold}。反応率は導出されない。

整数 \(2\)、\(8\)、\(20\)、\(28\)、\(50\)、\(82\)、\(126\) は席の数である。殻 \(N\) までの、スピンで重みを付けた振動子の閉包は \(C(N)=(N+1)(N+2)(N+3)/3\) である。\(N\ge 2\) では、さらに続く閉包 \(C(N)+2(N+2)\) は、その数に次の殻の最高 \(j\) 軌道の縮退を加えたものである。これらの整数は節の番号ではない。二重魔法核種は、各ローター類についてそのような数を一つ閉じる。層の幾何は整数を選ばず、安定の谷も定めない。

双対励起は、相対ラピディティとは別の座標である。励起 \(\psi\in(0,\pi)\) を固定すると、外側の井戸は吸引子をちょうど一つもち、それは臨界角 \(2\arctan(\gamma_s/\tau)\) が \(\psi\) に等しい位相にある。その臨界角は、無限遠と第一の節とのあいだで \(\pi\) から \(0\) まで厳密に落ち、\(\pi\) からの不足は \(2\ell/r\) から始まる。相対ラピディティは \(2\ell/\pi\) の外側で力を反転できない。双対励起は、あらゆる有限半径で反転できる。複合核がすでに存在するときの核分裂は、その反転である。核融合は、内向きのままである入口である。遅い中性子が核分裂を開きうるのは、まずラピディティの閾値を通過するときだけである。
\end{abstract}

\section{はじめに}
\label{sec:intro}

二重時空理論の核の節は、強い力を近接重力として読む。等スケール軌跡の上の重力チャンネルであり、そこでは干渉因子 \(\gamma_s\) が離散的な半径の集合で符号を変える~\cite{dst-main}。第二の結合は導入しない。同じ節は、その読みがしないことも明示している。ニュートンの組合せ \(G m_p^2\) は \(O(1)\) の核力強度ではない。節の近傍では \(1/\gamma_s\) が発散し、その発散は、すでに存在するチャンネルの運動学的増幅である。フェムトメートルと MeV の尺度は、長さ \(\ell\) と質量とを通して入り、\(\ell\) は導出されない。魔法数は共鳴条件として名づけられ、節からは取り出されない。核分裂と核融合は、層の遷移として名づけられる。

本稿は、それらの言明が指す層を計算する。第\ref{sec:force}節は、符号つきの力と、後の閾値がそれから組まれる二つの補助関数とを定める。第\ref{sec:layer}節は斥力反転そのものである。\(\gamma_s\) が負である最初の区間、壁の両側での力の向き、そして等スケールの力が間隔を外側の壁へ向けて駆動し、層を貫いては運ばない、その意味を述べる。第\ref{sec:rapidity}節は、相対ラピディティを通常セクターだけに置く。力が斥力になる閾値は位相の閉じた関数であり、静的な節と半径 \(2\ell/\pi\) とのあいだにだけ存在する。内向きの自由なラピディティは、静的な節より前にその閾値に出会う。捕獲、中性子誘起核分裂、核融合は、その入口を共有する。遅い中性子とは、それをまだ横切っていない中性子である。第\ref{sec:proton}節は、陽子対に電磁チャンネルを加える。二つのチャンネルは、長さが等しくないかぎり壁を共有せず、したがって陽子を複数もつ核は、引き合う対と斥け合う対とを同時にもつ。第\ref{sec:cold}節は、遅い荷電の入口である。ラピディティの閾値より下では、核の項はクーロン項に対抗し、閾値のすぐ下ではそれを上回る。閾値より上では、核の項はクーロン斥力に加わる。熱的なラピディティは、その始まりを静的な節に寄せる。第\ref{sec:magic}節は、席の数 \(2\)、\(8\)、\(20\)、\(28\)、\(50\)、\(82\)、\(126\) を節の列から分け、二重に閉じた核種を列挙する。第\ref{sec:fission}節は、間隔を固定した双対励起に戻る。その励起は、ラピディティと異なり、あらゆる半径で外側の井戸を反転できる。反転は、すでに形成された核の分裂である。第二の入り方ではない。

電磁の同伴稿は、同じ \(\gamma_s\) を電荷にすでに用いている~\cite{dst-shell}。以下の符号は、その力の重力への特殊化であり、粒子が陽子であるところでは電磁の項が加わる。離散の同伴稿は、核を陽子と中性子のローター類の順序つき積として表し、固定されたコンパクト化のそれぞれが、代数的に安定な質量数を有限個しかもたないことを証明する~\cite{dst-dio}。\(A\sim 300\) 付近の数値的な天井は、そこで既に述べられたヒューリスティックのままである。本稿の層の幾何の中には、それに代わるものはない。

\section{符号つきの近接力}
\label{sec:force}

一軸の上で、双対ローターの干渉因子は
\begin{equation}
\label{eq:gamma-ab}
\gamma_s(\alpha,\beta)
  =\cosh\frac{\alpha}{2}\cos\frac{\beta}{2}
  -\sinh\frac{\alpha}{2}\sin\frac{\beta}{2}
\end{equation}
である。等スケール軌跡は \(\alpha=\beta=\ell/r\) と置く。無次元の位相 \(x=\ell/(2r)\) では
\begin{equation}
\label{eq:gamma}
\gamma_s(x)=\cosh x\cos x-\sinh x\sin x,
\qquad
\gamma_s'(x)=-2\cosh x\sin x
\end{equation}
である。したがって \(\gamma_s(0)=1\) である。以下で用いるテイラー展開は
\begin{equation}
\label{eq:series}
\gamma_s(x)=1-x^{2}-\frac16 x^{4}+O(x^{6})
\end{equation}
である。閾値には、さらに二つの組合せが入る。それらは
\begin{align}
\label{eq:sigma}
\sigma(x)&=\cosh x\sin x-\sinh x\cos x,\\
\label{eq:tau}
\tau(x)&=\cosh x\sin x+\sinh x\cos x
\end{align}
である。どちらも原点で消える。導関数は \(\sigma'(x)=2\sinh x\sin x\) と \(\tau'(x)=2\cosh x\cos x\) である。和と差は
\begin{equation}
\label{eq:sumdiff}
\sigma+\gamma_s=e^{-x}(\sin x+\cos x),
\qquad
\sigma-\gamma_s=e^{x}(\sin x-\cos x)
\end{equation}
に崩れる。\((0,\pi)\) では \(\sigma(x)>0\) である。\(\sigma(0)=0\) であり、その区間で \(\sigma'\) は非負で、どの部分区間でも恒等的には零でないからである。\((0,\pi/2)\) では \(\tau(x)>0\) である。\(\tau(0)=0\) であり、そこでは \(\tau'>0\) だからである。

二つの核子質量の間隔を増す向きの力は
\begin{equation}
\label{eq:force}
F_r=-\frac{\kappa}{\gamma\,r^{2}}
\end{equation}
である。ここで \(\gamma\) は、対に割り当てたラピディティの上で評価した~\eqref{eq:gamma-ab} であり、\(\kappa>0\) は核のチャンネルの前因子である。\(\gamma>0\) ならば力は内向きである。\(\gamma<0\) ならば外向きである。等スケール軌跡の上では \(\gamma=\gamma_s(x)\) である。本文は同じ商を、ニュートンの前因子 \(G m_1 m_2\) で書いている~\cite{dst-main}。その同定は、核の井戸としては用いない。第\ref{sec:proton}節は、クーロンの前因子との食い違いの大きさを記す。半径と符号と二つの閾値は、\(\kappa\) に依らない。

\(\gamma\) の単純零点では、力は単純極をもつ。親理論の固有時因子は、その零点で正のままである。節は力の発散であり、時計の反転ではない~\cite{dst-main}。

\section{斥力反転層}
\label{sec:layer}

各区間 \([n\pi,(n+1)\pi]\) で、導関数~\eqref{eq:gamma} は \(-\sin x\) の符号をもつ。したがって \(\gamma_s\) はそこで狭義単調であり、零点をちょうど一つもつ。零点は \(\tan x=\coth x\) のところにある。有限な正の \(x\) では \(\coth x>1\) であるから、第 \(n\) 零点は
\[
n\pi+\frac{\pi}{4}<x_n<n\pi+\frac{\pi}{2}
\]
にある。最初の零点 \(x_1\) は \((\pi/4,\pi/2)\) における唯一の根である。数値的には \(x_1=0.93755\) であるから、
\[
\frac{r_1}{\ell}=\frac{1}{2x_1}=0.53330
\]
である。第二の零点は \(x_2=3.9274\) であり、\(r_2/\ell=0.12731\) である。それらの半径の比は \(r_2/r_1=x_1/x_2=0.23872\) である。本文と同じ限界
\[
\frac16<\frac{r_2}{r_1}<\frac{4}{5\pi}
\]
は、括弧 \(\pi/4<x_1<1\) と \(x_2\in(5\pi/4,3\pi/2)\) とから従う。上からの見積もりは \(x_1/x_2<1/(5\pi/4)=4/(5\pi)\) である。下からの見積もりは \(x_1/x_2>(\pi/4)/(3\pi/2)=1/6\) である。半径の大きい方から読むと、\(r>r_1\) では \(\gamma_s>0\)、\(r_2<r<r_1\) では \(\gamma_s<0\) であり、符号はこのあと交替する。区間 \(r_2<r<r_1\) が斥力反転層である。

その層の内部での \(\gamma_s\) の極値は \(\gamma_s(\pi)=-\cosh\pi\) であり、半径 \(\ell/(2\pi)\) にある。そこでの外向きの力は
\begin{equation}
\label{eq:plateau}
|F_r|
  =\frac{4\pi^{2}}{\cosh\pi}\,\frac{\kappa}{\ell^{2}}
  =3.4057\,\frac{\kappa}{\ell^{2}}
\end{equation}
である。より高い番号の平坦部は、より静かである。\(x=n\pi\) では同じ簡約が
\[
|F_r|=\frac{4\pi^{2} n^{2}}{\cosh(n\pi)}\,\frac{\kappa}{\ell^{2}}
\]
を与え、\(\cosh(n\pi)\) が \(e^{n\pi}/2\) のように増えるので、これは幾何的に落ちる。中心へ向かって限りなく増えるのは極の集合であり、極と極とのあいだの力ではない。

第一の節は、間隔の吸引子である。\(x\) が下から \(x_1\) に近づくと \(\gamma_s\downarrow 0\)、\(F_r\to-\infty\) となり、間隔は減り、位相は節へ向かって増える。\(x\) が上から \(x_1\) に近づくと、\(\gamma_s\) は負の値を通って \(\uparrow 0\) となり、\(F_r\to+\infty\) であるから、間隔は増し、位相は節へ戻る。同じ符号の組は、\(\gamma_s'\) が負である零点のすべて、すなわち各区間 \((2m\pi,(2m+1)\pi)\) で成り立つ。補区間では \(\gamma_s'\) が正であり、横断は負の値から正の値へ走り、零点の両側は間隔を零点から遠ざける。第二の節は、その第二の種類である。したがって反転層の内部では外向きの力が \(r_1\) を指し、\(r_2\) のすぐ内側では内向きの力が \(r_2\) から離れ、次の引力の室の奥へ向かう。等スケール軌跡の上では、大きく始まる間隔は第一の節へ向けて駆動され、層を貫いて \(r_2\) の向こうの室へは入らない。落下のラピディティが増えることを許せば、経路はその軌跡を離れる。第\ref{sec:rapidity}節の閾値に、\(x_1\) より厳密に小さい位相で出会い、奥の室には達しない。

節の外側での~\eqref{eq:force} の動径積分は、下に有界でないまま落ちる。電磁の同伴稿がクーロンの井戸に用いるのと同じ置換による~\cite{dst-shell}。無限遠からの経路は極を横切らず、上端の位相が \(x_1\) に近づくにつれて積分は発散する。反転層はその極の向こう側にある。局所的な力は \(\gamma_s\neq 0\) のところで定義され、外側の井戸の積分は節を越えて続かない。

長さは荷電パイオンのコンプトン波長 \(\ell=\lambda_\pi\approx 1.41\,\mathrm{fm}\) である。本文は、そのときすべての等スケール節が \(1\,\mathrm{fm}\) 未満にあり、最外の節が区間 \(0.44\)--\(0.91\,\mathrm{fm}\) に限られることを記している~\cite{dst-main}。中心値は \(r_1=0.752\,\mathrm{fm}\) である。同定は \(\ell\) の導出ではない。反転層を \(1\)--\(2\,\mathrm{fm}\) に置くには、より長い長さが要る。節の半径の比は、長さに依らない。

核子間相互作用の符号は、この長さを一つの窓に置く。散乱データに合わせた現実的なポテンシャルは、\(\tfrac12\,\mathrm{fm}\) までの間隔で斥力であり、\(1\,\mathrm{fm}\) から \(2\,\mathrm{fm}\) まで引力である~\cite{wiringa1995}。間隔 \(\tfrac25\,\mathrm{fm}\) はその芯の内側にある。遅い対では、そこでの斥力はその間隔を第一の節の内側に置く。節は \(14/15\) を厳密に上回るので、
\[
\ell>\frac{28}{15}\cdot\frac{2}{5}=\frac{56}{75}\,\mathrm{fm}
\]
である。\(1\,\mathrm{fm}\) での引力は、その間隔を外側の井戸に置く。節は \(19/20\) を厳密に下回るので、\(\ell<(19/10)\,\mathrm{fm}\) である。二つの半径は二倍だけ離れている。節を五分の一広げる移動は、\(100\,\mathrm{MeV}\) で既に得られている五パーセントより大きく、その移動のあとでも \(\tfrac25\,\mathrm{fm}\) と \(\tfrac12\,\mathrm{fm}\) はともに \(1\,\mathrm{fm}\) の内側に残る。荷電パイオンの波長は、\(56/75\,\mathrm{fm}\) と \(19/10\,\mathrm{fm}\) のあいだに厳密にある。

この長さのもとで、第一の節は \(14/19\) と \(213/280\,\mathrm{fm}\) のあいだに厳密にあり、したがって \(\tfrac12\,\mathrm{fm}\) を厳密に上回り、\(1\,\mathrm{fm}\) を厳密に下回る。第二の節は \(\tfrac12\,\mathrm{fm}\) を厳密に下回るので、芯の尺度は反転層の中にある。ラピディティの壁 \(2\ell/\pi\) は \(8/9\) と \(142/157\,\mathrm{fm}\) のあいだに厳密にあり、したがって \(1\,\mathrm{fm}\) を厳密に下回る。有限なラピディティは、引力区間の内側の縁で力を反転しない。

測定された電荷半径は、電荷分布の大きさである。対の間隔ではない。表~\ref{tab:radii} は、アンジェリとマリノヴァの二乗平均平方根と、引用された不確かさとを記す~\cite{angeli2013}。各下端、すなわち中心値からその不確かさを引いた値は、\(1\,\mathrm{fm}\) を上回る。したがってこの長さのもとでは、それぞれがラピディティの壁の外側にある。

\begin{table}[htbp]
\centering
\caption{電荷の二乗平均平方根半径。電荷半径は対の間隔ではない。各下端は、\(\ell=\lambda_\pi\) のもとでラピディティの壁の外側にある。}
\label{tab:radii}
\begin{tabular}{lcc}
\hline
核種 & \(R\) (fm) & 不確かさ (fm) \\
\hline
\({}^{2}\mathrm{H}\) & \(2.1421\) & \(0.0088\) \\
\({}^{4}\mathrm{He}\) & \(1.6755\) & \(0.0028\) \\
\({}^{16}\mathrm{O}\) & \(2.6991\) & \(0.0052\) \\
\({}^{40}\mathrm{Ca}\) & \(3.4776\) & \(0.0019\) \\
\({}^{48}\mathrm{Ca}\) & \(3.4771\) & \(0.0020\) \\
\({}^{208}\mathrm{Pb}\) & \(5.5012\) & \(0.0013\) \\
\hline
\end{tabular}
\end{table}

飽和密度 \(n_0=0.16\,\mathrm{fm}^{-3}\) における核子あたりの平均体積は \(25/4\,\mathrm{fm}^{3}\) である。その立方根は \(46/25\) と \(37/20\,\mathrm{fm}\) のあいだに厳密にある。その間隔はラピディティの壁の外側にある。それが与える上界は \(\ell<(19/10)\cdot(37/20)\,\mathrm{fm}\) であり、\(1\,\mathrm{fm}\) での引力が与える上界より弱い。

親理論は、共有された席の排他を、別の統計的要請としてではなく、この層として読む。重なる同種フェルミオンは、第一の節の外向きの側に出会う~\cite{dst-main}。その読みを、ここでは核子に採用する。これはスピン統計定理の導出ではない。陽子と中性子は異なるローター類である~\cite{dst-dio}。したがって同じ文は、それだけでは、一つの間隔にある陽子--中性子対を禁じない。

束縛の尺度は幾何の中にはない。本文は、鉄付近の \(\mathrm{BE}/A\sim 8\,\mathrm{MeV}\) と、\(40\)--\(50\,\mathrm{MeV}\) の程度の光学模型の井戸深さとを、別々の入力として引用する。どちらも入力のままである。密度の飽和、すなわち \(R=r_0 A^{1/3}\) という形は、層の描像が主張する規則性であり、本文はすでに、\(r_0\) の二つのありふれた数値規約が一割の水準で食い違うことを記している。本稿はそのあいだを選ばない。

\section{相対ラピディティ}
\label{sec:rapidity}

相対運動は、通常セクターのブーストである。親理論では、そのブーストのラピディティは \(\tanh\varphi=|v|/c\) を満たす正のパラメーター \(\varphi\) であり、ラピディティ成分の平方和の平方根である~\cite{dst-main}。ここで採用する読みは、そのパラメーターを近接の値に加え、双対角は近接の値に残す。
\begin{equation}
\label{eq:ansatz}
\alpha=\frac{\ell}{r}+\varphi,
\qquad
\beta=\frac{\ell}{r},
\qquad
\varphi\ge 0.
\end{equation}
\(\varphi\neq 0\) となった瞬間に、等スケールは破れる。対は、不整合スカラー \(J\) が消える軌跡を離れる。力の符号は、なお~\eqref{eq:gamma-ab} の符号である。逆ブースト \(\alpha\mapsto-\alpha\) は別の元である。\(\sinh(\alpha/2)\) の前の符号を変え、干渉因子を上げる。捕獲は正のラピディティの上で読む。それはローターが実際に担うパラメーターである。

\(\delta=\varphi/2\) と書く。~\eqref{eq:gamma-ab} へ代入すると
\begin{equation}
\label{eq:gamma-phi}
\gamma(x,\delta)
  =\cosh\delta\,\gamma_s(x)-\sinh\delta\,\sigma(x)
\end{equation}
を得る。閉じた形
\begin{equation}
\label{eq:gamma-closed}
\gamma(x,\delta)
  =\frac12 e^{\delta+x}(\cos x-\sin x)
  +\frac12 e^{-(\delta+x)}(\cos x+\sin x)
\end{equation}
は~\eqref{eq:sumdiff} から従い、符号を明瞭にする。\(0\le x\le\pi/4\) では両方の係数が非負であるから、あらゆる有限な \(\delta\) で \(\gamma(x,\delta)>0\) である。半径 \(r\ge 2\ell/\pi\) では、どの相対ラピディティも力を反転しない。

\(\pi/4<x<x_1\) では \(0<\gamma_s<\sigma\) であるから、比 \(\rho=\gamma_s/\sigma\) は \((0,1)\) にある。因子~\eqref{eq:gamma-phi} は
\[
\tanh\delta=\rho(x)
\]
で消え、したがって
\begin{equation}
\label{eq:phi-star}
\varphi_*(x)
  =2\operatorname{artanh}\rho(x)
  =\ln\frac{\sin x+\cos x}{\sin x-\cos x}-2x
\end{equation}
で消える。この閾値より下では \(\gamma>0\) であり、力は内向きである。閾値より上では \(\gamma<0\) であり、力は外向きである。\(x\downarrow\pi/4\) では分母 \(\sin x-\cos x\) が消え、\(\varphi_*\to\infty\) となる。\(x\uparrow x_1\) では \(\gamma_s\to 0\) となり、\(\varphi_*\to 0\) となる。

閾値は位相とともに厳密に落ちる。\((0,\pi)\) の上で
\begin{equation}
\label{eq:rho-derivative}
\gamma_s'\sigma-\gamma_s\sigma'
  =-2\sin^{2}x
\end{equation}
である。この恒等式は、\(\gamma_s'=-2\cosh x\sin x\) と \(\sigma'=2\sinh x\sin x\) の代入である。残る括弧は \(\cosh x\,\sigma+\sinh x\,\gamma_s\) であり、その括弧は \(\sin x\) に等しい。したがって \(\sin x\neq 0\) のところでは \((\gamma_s/\sigma)'<0\) である。とくに \(\rho\) は \((0,x_1)\) の上で \(+\infty\) から \(0\) まで落ち、ちょうど \(\pi/4\) で \(1\) に等しく、\(\varphi_*\) は \((\pi/4,x_1)\) の上で \(+\infty\) から \(0\) まで落ちる。与えられたラピディティが閾値の上に座る半径 \(r_*=\ell/(2x_*)\) は、したがって \(\varphi\) とともに上がる。\(\varphi=0\) では \(r_1\) から始まり、\(\varphi\to\infty\) で \(2\ell/\pi\) に近づく。極限の比は
\begin{equation}
\label{eq:limit-ratio}
\frac{2\ell/\pi}{r_1}=\frac{4x_1}{\pi}=1.1937
\end{equation}
である。反転層が外へ厚くなる分は、どのラピディティでも五分の一より小さい。

窓の内部では、~\eqref{eq:gamma-phi} の \(\delta\) に関する導関数は厳密に負である。より高いラピディティは \(\gamma\) を下げる。\(\gamma\) がなお正であるあいだ、内向きの力は強まる。\(\varphi_*\) では内向きの力が発散する。\(\varphi_*\) の直後では外向きの力が発散し、\(\varphi\) をさらに増すと \(\gamma\) はより負へ駆動されるので、外向きの大きさは無限大から、~\eqref{eq:gamma-closed} の第一項が読む超相対論的な大きさへ向かって落ちる。高い相対速度は、固定された斥力の法則を掛けるのではない。符号の変わる場所を動かす。窓の中の与えられた間隔の斥力領域は \(\varphi>\varphi_*(x)\) という領域であり、その領域の入口で斥力は任意に大きい。

窓の外側では、同じ導関数が一度だけ消え、それは \(\tanh\delta=\sigma/\gamma_s<1\) のところである。そこでは \(\gamma\) が正の極小をもつ。小さなラピディティは引力を強め、大きなラピディティはそれを零へ向かって弱める。~\eqref{eq:gamma-closed} が \(e^{\varphi/2}\) のように増えるからである。\(r>2\ell/\pi\) での力は、あらゆる有限の速さで内向きのままである。

小さな \(\varphi\) では、閾値は静的な節の近くにある。節の条件は \(\gamma_s(x_1)=0\) であり、
\[
\gamma_s'(x_1)=-2\cosh x_1\sin x_1
\]
である。\(\gamma_s=0\) の点での~\eqref{eq:sumdiff} は \(\sigma(x_1)=\sin x_1/\cosh x_1\) を与える。節における \(\varphi_*\) の導関数はちょうど \(-4\cosh^{2}x_1\) であるから、
\begin{equation}
\label{eq:skin}
x_1-x_*
  \sim
  \frac{\varphi}{4\cosh^{2}x_1},
\qquad
\frac{1}{4\cosh^{2}x_1}=0.11528
\end{equation}
である。一次まででは、半径の分数的な広がりはこの係数を \(x_1\) で割ったものであり、\(0.12295\,\varphi\) である。位相の間隙そのものは、あらゆる正のラピディティにおいて \(\varphi/(4\cosh^{2}x_{1})\) より厳密に小さく、分数的な広がり \((r_{*}-r_{1})/r_{1}\) は \(\varphi/\pi\) より厳密に小さい。遅い中性子の斥力の表皮は、\(r_1\) のその分数的な厚さの近傍である。

実験室の運動エネルギーから \(\varphi\) への写像は、中性子の静止エネルギー \(m_n c^{2}=939.6\,\mathrm{MeV}\) を入力として用い、\(\cosh\varphi=1+T/(m_n c^{2})\) を通る。そのラピディティでの閾値半径は、\(\varphi_*(\ell/(2r_*))=\varphi\) を満たす唯一の \(r_*\) である。

\begin{table}[htbp]
\centering
\caption{運動エネルギー \(T\) の中性子の反転半径。極限の比は \(4x_1/\pi=1.1937\) である。}
\label{tab:neutron}
\begin{tabular}{lcc}
\hline
\(T\) & \(\varphi\) & \(r_*/r_1\) \\
\hline
\(1\,\mathrm{eV}\) & \(4.6\times 10^{-5}\) & \(1.00001\) \\
\(1\,\mathrm{keV}\) & \(0.00146\) & \(1.00018\) \\
\(1\,\mathrm{MeV}\) & \(0.04613\) & \(1.00563\) \\
\(14\,\mathrm{MeV}\) & \(0.1724\) & \(1.0206\) \\
\(100\,\mathrm{MeV}\) & \(0.4574\) & \(1.0517\) \\
\(1\,\mathrm{GeV}\) & \(1.353\) & \(1.124\) \\
\hline
\end{tabular}
\end{table}

熱中性子は \(\varphi\) が \(10^{-5}\) の程度であり、表の精度では静的な節と区別できない。その表皮を除く窓の全体で、それに働く力は内向きであり、表皮に近づくにつれて内向きの力は大きくなる。\(14\,\mathrm{MeV}\) の中性子が閾値に出会うのは、静的な節の外、\(r_1\) の二パーセントだけである。そのエネルギーは、重水素--三重水素の核融合が放出する中性子のエネルギーである。核融合を誘起する中性子のエネルギーではなく、それを放出する反応は荷電の入口であり、以下で扱う。\(100\,\mathrm{MeV}\) でのずれは五パーセントである。幾何は、数 MeV の中性子をハードコア半径の二倍で引き返させない。斥力の表皮を \(\varphi\) に比例して厚くし、厳密な壁 \(2\ell/\pi\) までである。

\(\ell=\lambda_\pi\approx 1.41\,\mathrm{fm}\) のもとでは、その壁は \(2\ell/\pi=0.898\,\mathrm{fm}\) にある。境界 \(7/5<\lambda_\pi<71/50\) から読むと、同じ壁は \(8/9\) と \(142/157\,\mathrm{fm}\) のあいだに厳密にある。有限なラピディティは、斥力の始まりをその壁の向こうへは動かさない。静的な節は \(0.752\,\mathrm{fm}\) にあり、同じ境界はその節を \(14/19\) と \(213/280\,\mathrm{fm}\) のあいだに厳密に置く。壁と節との比は、長さに依らない。

瞬間の判定は、まだ捕獲断面積ではない。力が内向きのままであり、中性子がなお内へ動いているあいだ、ラピディティ自身が増える。相対運動量を \(\mu c\sinh\varphi\)、\(\mu\) を換算質量として、
\begin{equation}
\label{eq:infall}
\mu c\cosh\varphi\,\frac{d\varphi}{dt}
  =\frac{\kappa}{\gamma r^{2}},
\qquad
\frac{dr}{dt}=-c\tanh\varphi
\end{equation}
が、\(\gamma>0\) であるかぎり成り立つ。同値に、
\[
\sinh\varphi\,d\varphi
  =-\frac{\kappa}{\mu c^{2}}\,\frac{dr}{\gamma r^{2}}
\]
である。右辺は \(\gamma\) を通して \(\varphi\) に依る。\(\gamma\) が正であるあいだ、間隔が落ちるにつれてラピディティは増え、窓の上では閾値が落ちるので、余裕 \(\varphi_*-\varphi\) は両側から狭まる。それは \(x_1\) より厳密に小さい位相で閉じる。閾値より下の値から増える非負のラピディティは、節より前に、閾値に一度出会う。チャンネルがなお引力であるあいだ、静的な節には達しない。引力の側で節に達するには、ラピディティが下がらなければならない。放射、励起、光学ポテンシャルの虚部は~\eqref{eq:force} の中になく、本稿はそれらを導出しない。したがって~\eqref{eq:infall} を複合核の中まで積分せず、捕獲断面積が \(1/v\) に従って落ちるという実験室の規則も導出しない。符号つきの力が定めるのは、力の符号が変わる面と、内向きの自由なラピディティがその面に静的な節より前に出会うという事実である。間隔 \(r\in(r_1,2\ell/\pi)\) に \(\varphi>\varphi_*(\ell/(2r))\) で到着する中性子は、すでに極の外向きの側にいる。同じ間隔への、より遅い到着は内向きの側にあり、そのあとラピディティが増え続ければ、\(r_1\) より前に極に出会う。

遅い中性子への選好は入口に属し、入口は共有される。捕獲は中性子を受け入れ、内向きのままである。中性子誘起核分裂は同じ中性子を受け入れ、そのあと分裂しうる。核融合は相手の核を受け入れ、内向きのままである。\(\varphi_*\) より上ではチャンネルはすでに外向きであるから、三つとも始まらない。式が与えるのは、その共通の拒否である。相手が内に入ったあと、どの反応が続くかは選ばない。

荷電の入射粒子は、この核の条件を保ち、次の節のクーロン項を加える。内向きの核の項がそのクーロン項を上回るかどうかは、第\ref{sec:cold}節の荷電の入口である。入射粒子は電荷を担うので、その入口は遅い中性子の反応ではない。

\section{カオス的な陽子力}
\label{sec:proton}

陽子は両方のチャンネルを担う。重力チャンネルは~\eqref{eq:force} である。同符号の対の電磁チャンネルは、それ自身の長さ \(\ell_E\) をもつそれ自身の等スケール軌跡の上で、同伴稿の符号つきクーロン力を寄与する~\cite{dst-shell}。二つの陽子の間隔を増す力は
\begin{equation}
\label{eq:two-channel}
F_r
  =-\frac{\kappa}{\gamma(x,\varphi/2)\,r^{2}}
  +\frac{k e^{2}}{\gamma_s(x_E)\,r^{2}},
\qquad
x=\frac{\ell}{2r},
\quad
x_E=\frac{\ell_E}{2r},
\end{equation}
であり、\(k=1/(4\pi\epsilon_0)\) である。ラピディティは慣性の対に割り当てる。電荷の対は等スケール因子を保つ。中性子は \(e=0\) の場合である。

\paragraph{長さの選び方。}\label{rule:length}
各項の長さは、それを担う粒子によって選ぶ。電磁の長さ \(\ell_E\) は、孤立した電荷としてふるまう電荷、すなわち電子と自由な陽子に属する。重力の長さ \(\lambda_\pi\) は、質量をもつ粒子、主に陽子と中性子、わずかに電子の相互作用に属する。原子核の中の陽子は本節の陽子のカオスにあり、孤立した電荷としてはふるまわない。したがって原子核の水準での相互作用は、そのクーロン項も含めて \(\lambda_\pi\) で読む。

二つの長さは独立な入力である。\(\lambda=\ell_E/\ell\) と書く。電磁の節は、重力の位相で読むと \(x_n/\lambda\) にある。それらがあらゆる \(n\) で重力の節と一致するのは、\(\lambda=1\) のときに限る。その特別な場合、静止した~\eqref{eq:two-channel} は、前因子 \(k e^{2}-\kappa\) をもつ単一の層の列である。\(\kappa>k e^{2}\) ならば正味の力は重力の符号をもつ。\(\kappa<k e^{2}\) ならばクーロンの符号をもつ。壁の第二の集合はない。長さの規則により、組み上がった原子核の陽子の対はこの場合で読む。このとき室は単一の列の層であり、内向きと外向きを交互にとり、以下のカオス的な陽子力の符号胞の読みはそのまま当てはまる。\(\lambda\neq 1\) の互い違いは、陽子の一方がなお孤立した電荷としてふるまう対を表す。

\(\lambda\neq 1\) のとき、壁は互い違いになる。どちらの因子の節も含まない各開区間が一つの室である。両方の分母が符号を保ち、各項の向きは一定である。有界な位相区間に出会う室の数は、内側の端点を原点へ押し込むにつれて限りなく増える。各チャンネルが、長さ \(\pi\) の区間ごとに節を一つ寄与するからである。\(\lambda\) が無理数ならば、重力の節は電磁の節と一致せず、どの壁も一方の項だけの単純極であり、位相における室の長さの列は周期的でない。大きな間隔から内へ向かうときに最初に出会う壁は、\(\lambda<1\) ならば重力の節、\(\lambda>1\) ならば電磁の節である。より短い到達距離の反転の外側にある有限なクーロン斥力は、この二つの順序のうちの第一である。\(\ell_E<\ell\) であり、電荷の因子は \(r_1\) でなお正であり、重力の極が \(r_1\) のある近傍でその有限な項を追い越す。近傍の幅は \(\kappa/(k e^{2})\) が定め、純数ではない。逆の順序は、重力の節の外側にクーロンの極を置き、無限遠からの古典的な接近はそこで引き返される。

壁の安定性は、その項の前の符号と \(\gamma_s'\) の符号とで定まる。\(\gamma_s\) が零を下向きに横切る零点を考える。したがって \(\gamma_s'<0\) である。重力の側では、隣り合う両方の室が間隔を零点へ向けて駆動する。第\ref{sec:layer}節のとおりである。電磁の側では \(1/\gamma_s\) の前因子が負ではなく正であり、隣り合う両方の室が間隔を零点から遠ざける。すぐ外側では同符号の電荷が斥け合い、間隔は増える。すぐ内側では \(\gamma_s<0\) であり同符号の電荷が引き合うので、間隔は減り、位相は零点から離れる。\(\gamma_s\) が零を上向きに横切る零点は、各チャンネルで逆の安定性をもつ。したがって一つのチャンネルの隣り合う壁は、間隔の吸引子と反発子とのあいだで交替する。最外の重力の壁は吸引子である。最外の電磁の壁は反発子である。

互い違いは、比の確定した値ではない図示 \(\lambda=2\) で、すでに見える。重力の位相 \(x\) が零から増えると、最初の壁は \(x=x_1/2=0.46878\) にある電磁の節である。それを越え、重力の節 \(x_1=0.93755\) までのあいだ、電磁の因子は負であるからクーロン項は内向きであり、重力の項もなお内向きである。その室は両方のチャンネルから引く。\(x_1\) を越え、次の電磁の節 \(x=x_2/2=1.9637\) までのあいだ、重力の項は外向きになり、クーロン項は内向きのままである。正味の向きは、そのとき比較
\begin{equation}
\label{eq:contest}
\frac{\kappa}{|\gamma|}
  \gtrless
  \frac{k e^{2}}{|\gamma_E|}
\end{equation}
であり、これは入力 \(\kappa\) と二つの分母とに依る。それより後の各壁は、二つの項のうちちょうど一つを反転する。この列を内へ動く陽子対は、二重に内向きの室と、争われる室と、クーロン項が再び反転した室とのあいだで交替するであろう。古典的には、対は吸引する壁を自由に横切れない。力はそこで発散する。交替は室の構造であり、極小を一つもつ滑らかな井戸ではない。

核の中では、対は一つの間隔を共有しない。一直線上の三つの陽子を、\(0\)、\(s\)、\(s+t\) に考え、\(s\)、\(t\)、\(s+t\) はどの節の上にもないとする。\(F(s)\) を、間隔 \(s\) における~\eqref{eq:two-channel} とし、その間隔を増すときに正とする。中央の陽子に働く右向きの力は \(F(s)-F(t)\) である。左の間隔が内向きの室にあり \(F(s)<0\)、右の間隔が外向きの室にあり \(F(t)>0\) ならば、差は負である。左の相手は中央の陽子を左へ引き、右の相手はそれを左へ押す。二つの力は同じ向きを指す。それらの間隔の一つを節の上へ運ぶ変位は、対応する \(F\) を無限大を通して反転させる。中央の陽子に働く力の変化は、壁を横切る任意に小さい変位において非有界である。壁の両側の配置は、力において近くない。

これがカオス的な陽子力である。動径の二体問題がカオス力学系であるという主張ではない。一つの間隔と、定められた \(F(r)\) とは、自由度が一つである。ここでいうカオスは配置である。いくつかの対がいくつかの室に同時に座り、あるものは内向き、あるものは外向きであり、各陽子の上のベクトル和が滑らかなのは符号胞の内部だけである。胞の境界は極である。核は、一つの引力の井戸の中の対の集まりではない。対の間隔の符号胞であり、どの対も壁に出会わないあいだだけ安定で、一つの対を壁の上へ運ぶ変位に敏感である。

ニュートンの同定は、壁から離れたところでの争い~\eqref{eq:contest} を極端にする。CODATA の値は
\begin{equation}
\label{eq:ratio}
\frac{G m_p^{2}}{k e^{2}}=8.094\times 10^{-37},
\qquad
\frac{k e^{2}}{G m_p^{2}}=1.236\times 10^{36}
\end{equation}
を与える~\cite{codata2018}。\(\kappa\) を \(G m_p^2\) で置き換えると、重力の項がクーロン項を上回るのは、\(|\gamma|\) が \(|\gamma_E|\) よりその因子 \(1.236\times 10^{36}\) だけ小さいところだけである。平坦部では両方の因子が一の程度かそれより大きく、符号を決めるのはクーロン項である。本文はすでに、\(G m_p^{2}/\hbar c\sim 10^{-38}\) を \(O(1)\) の核力強度と混同することを拒んでいる~\cite{dst-main}。その拒否は保たれる。核のチャンネルの前因子 \(\kappa\) は入力である。室の壁、その安定性、ラピディティの閾値は、それに依らない。争われる室の内部でどちらの項が勝つかは、依る。

陽子--中性子対は重力の項をもち、クーロン項をもたない。陽子--陽子対は両方をもつ。外側の重力の井戸は前者を束縛でき、\(k e^{2}/\gamma_E\) が \(\kappa/\gamma\) を上回れば後者の束縛に失敗する。両方の因子が正であるところでは、これが~\eqref{eq:contest} の争われない向きである。実験室の重陽子は陽子--中性子の束縛状態であり、二重陽子は非束縛である。本稿は \(2.2\,\mathrm{MeV}\) の束縛を計算せず、中性子--中性子散乱長も導出しない。節の外側の間隔にある二つの中性子は、共有された席を排他する極の上にはいない。したがって \(\gamma_s\) の符号だけでは、それらが非束縛であると決まらない。

\section{遅い荷電の入口}
\label{sec:cold}

クーロン項は、それだけでは、荷電の対が入れるかどうかを決めない。\(\rho=k e^{2}/\kappa\) と書き、間隔を増す力を、両方の分母が零でないところでは
\begin{equation}
\label{eq:charged-net}
F_r
  =\frac{\kappa}{r^{2}}
  \left(-\frac{1}{\gamma}+\frac{\rho}{\gamma_E}\right)
\end{equation}
と書く。力の符号は、第二の因子の符号である。本節を通して、電磁の因子 \(\gamma_E\) は正とする。これは、上で既に用いた順序 \(\ell_E<\ell\) である。電荷の長さの方が短く、大きな間隔から最初に出会う壁は重力の節であり、その節ではクーロン項はなお有限である。この順序は本節の仮定である。第~\ref{rule:length} 節の長さの規則により、原子核の水準での対は \(\ell_E=\ell\) で読む。

窓 \(\pi/4<x<x_1\) の上では、干渉因子は半ラピディティ \(\delta=\varphi/2\) が上がるにつれて厳密に落ちる。閾値より上では \(\gamma\) は負であるから、\(-1/\gamma\) は正であり、核の項はクーロン項に加わる。正味の係数は、そのとき \(\rho/\gamma_E\) より厳密に大きい。斥力はクーロン斥力を上回り、閾値の直後では任意に大きい。クーロン障壁を横切るために上げられたラピディティは、窓の内部で、障壁に加わった核の斥力に出会う。

閾値より下では \(\gamma\) は正であり、閾値で零へ落ちるので、\(-1/\gamma\) は下限なしに落ちる。あらゆる有限で正の \(\rho\) と、あらゆる正の \(\gamma_E\) とに対して、閾値の下側に接するラピディティの区間があり、その上で正味の係数は負である。そこでは核の項がクーロン項を上回る。

閾値の位相は、ラピディティが落ちるにつれて節へ向かって動く。したがって、より小さい正のラピディティは、より小さい半径まで内向きのままである。位相の間隙と分数的な広がりとは、遅い中性子の表皮について上で既に定めたものである。\(\cosh\varphi=1+T/(mc^{2})\) であり、
\[
0<\frac{T}{mc^{2}}\le 10^{-10}
\]
であるときには、その広がりは \(5\times 10^{-6}\) より厳密に小さい。室温の核子または重陽子は、この種のものである。熱エネルギーは電子ボルトの数百分の一であり、静止エネルギーは GeV の程度であるから、その比は \(10^{-10}\) より下にある。斥力の表皮は、そのとき、表~\ref{tab:neutron} の精度ではすでに見えない、\(1\,\mathrm{eV}\) の中性子の表皮よりも薄い。そのようなラピディティでは、核の斥力の始まりは静的な節と区別できず、核の項がクーロン項を上回る近傍は、その節に接している。

半径 \(2\ell/\pi\) の外側では、核の項は、いかなる有限なラピディティでも符号を変えない。ラピディティが増すにつれて干渉因子は限りなく増し、核の引力は零へ向かい、あるラピディティを越えれば正味の力は外向きである。対はクーロン斥力に残される。その半径の外側の間隔で速度を落とすことは、速い通過が消した核の引力を戻す。その側で引力が最も強いラピディティは、すでに位置の定められたもの、すなわちラピディティについての導関数が消えるところのままである。

低温を有利にするのは、極の側と、極の位置である。閾値より下では、ラピディティが上がるほど内向きの係数はより負になるので、有利さは静止におけるより強い引き込みではない。

第\ref{sec:rapidity}節の自由な落下は、なおラピディティを上げ、上がるラピディティはなお節より前に極に出会う。そのあと自由に落とされた遅い荷電の対は、その極の上で引き返される。低温の入口とは、間隔が節へ近づくあいだ、ラピディティが落ちていく閾値より下に保たれる入口である。正のイオンの正八面体の軸の上では、電磁チャンネルが同符号の電荷二つを、第一の電荷の節の外側の戻す席で保持し、その席は間隔を、電荷の節を貫く自由な電磁の落下から外す~\cite{dst-cold}。電子の換算コンプトン波長のもとで、その席の籠の半径 \(2\ell\) は \(772.2\,\mathrm{fm}\) と \(772.4\,\mathrm{fm}\) のあいだに厳密にあり、パラジウムの正八面体空隙は約 \(2\,\text{\AA}\) で、約 \(250\) 倍大きい。力は、反応率も、熱の過剰も、冷たい荷電の入口が起きるかどうかを問うてきた実験のあいだの選択も、導出しない。それが定めるのは符号である。閾値より下に留まる経路の上では、\(\gamma_E\) が正であるかぎり、あらゆる有限な \(\rho\) に対して、正味の力は閾値で終わるある区間の全体で内向きである。二つの長さの逆の順序は、その区間に達するより前に、クーロンの極で古典的な接近を引き返させる。

\section{魔法数と核種}
\label{sec:magic}

節の半径は \(\Theta(1/n)\) のように尺度をもつ。その列は、整数 \(2\)、\(8\)、\(20\)、\(28\)、\(50\)、\(82\)、\(126\) を含まない。それらの整数は席を数える。電磁の同伴稿が \(2\)、\(8\)、\(18\)、\(32\) を、相次ぐ角の次数の上の掌性の席二つとして読み、節のラベルとして読むことを拒むのと同じ意味である~\cite{dst-shell}。核の数は別の算術である。経験的な閉包が振動子の殻を分裂させるからである。その分裂は \(\gamma_s\) からは導出されない。以下に記すのは数そのものであり、反転層の性質と取り違えないためである。

三次元の等方振動子の殻 \(N=0,1,2,\dots\) は、軌道の縮退 \((N+1)(N+2)/2\) をもつ。各軌道にスピンの席を二つ置くと、状態は \((N+1)(N+2)\) 個になる。殻 \(N\) までの累積は
\begin{equation}
\label{eq:C}
C(N)
  =\sum_{k=0}^{N}(k+1)(k+2)
  =\frac{(N+1)(N+2)(N+3)}{3}
\end{equation}
である。最初の値は \(C(0)=2\)、\(C(1)=8\)、\(C(2)=20\)、\(C(3)=40\)、\(C(4)=70\)、\(C(5)=112\) である。殻 \(N+1\) の最高の軌道角運動量は \(l=N+1\) であり、状態 \(j=l+1/2\) の縮退は \(2j+1=2(N+2)\) である。その一つの軌道を累積に加えると、\(N\ge 2\) に対して
\begin{equation}
\label{eq:magic}
C(N)+2(N+2)
\end{equation}
を得る。すなわち \(28\)、\(50\)、\(82\)、\(126\) である。観測される閉包は \(C(0)\)、\(C(1)\)、\(C(2)\)、および~\eqref{eq:magic} である。\(N\ge 3\) の純振動子の閉包 \(C(N)\) は、それらの中にない。この算術は殻模型の算術である。次の殻の高い \(j\) の軌道の一つを、その下の閉じた殻とともに数える。層の幾何は、その軌道が、別の軌道ではなく、動く方である理由を言わず、本稿も欠けた理由を補わない。荷電独立性とスピン--軌道結合は、本文におけるとおり、節へ還元されていない双対ローターの共鳴条件のままである。

陽子の席と中性子の席は、二つのローター類の席である~\cite{dst-dio}。各類は上の数をもつ。電磁の項と、第\ref{sec:proton}節のカオス的な符号胞とは、陽子の類に働く。中性子の類には働かない。したがって整数は一致しえても、束縛は一致しない。配置に陽子を一つ加えることは、すでに存在する陽子対にわたるクーロン和を、\(\gamma_E>0\) である各室に加え、符号胞に陽子対を一つさらに加える。中性子を加えても、どちらも加わらない。重い核種の安定な組成が中性子過剰へ曲がるのは、そのためである。半経験的質量公式における \(Z^{2}/A^{1/3}\) の係数は、その曲がりによっては定まらず、ここでは導出されない。

二重魔法核種とは、陽子数と中性子数のそれぞれが \(2\)、\(8\)、\(20\)、\(28\)、\(50\)、\(82\)、\(126\) の一つである核種である。

\begin{table}[htbp]
\centering
\caption{二重に閉じた席数。二つの数の閉包は、崩壊に対する安定性ではない。}
\label{tab:doubly}
\begin{tabular}{lccc}
\hline
核種 & \(Z\) & \(N\) & \(A\) \\
\hline
\({}^{4}\mathrm{He}\) & \(2\) & \(2\) & \(4\) \\
\({}^{16}\mathrm{O}\) & \(8\) & \(8\) & \(16\) \\
\({}^{40}\mathrm{Ca}\) & \(20\) & \(20\) & \(40\) \\
\({}^{48}\mathrm{Ca}\) & \(20\) & \(28\) & \(48\) \\
\({}^{56}\mathrm{Ni}\) & \(28\) & \(28\) & \(56\) \\
\({}^{78}\mathrm{Ni}\) & \(28\) & \(50\) & \(78\) \\
\({}^{100}\mathrm{Sn}\) & \(50\) & \(50\) & \(100\) \\
\({}^{132}\mathrm{Sn}\) & \(50\) & \(82\) & \(132\) \\
\({}^{208}\mathrm{Pb}\) & \(82\) & \(126\) & \(208\) \\
\hline
\end{tabular}
\end{table}

これらのうち、\({}^{4}\mathrm{He}\)、\({}^{16}\mathrm{O}\)、\({}^{40}\mathrm{Ca}\)、\({}^{48}\mathrm{Ca}\)、\({}^{208}\mathrm{Pb}\) は安定である。他は谷の外にある。席の数は、どちらの場合も閉じる。\(\mathrm{BE}/A\) の最大は \(A=62\) の近く、閉包 \(28\) と \(50\) とのあいだにあり、\({}^{56}\mathrm{Ni}\) の上にはない。陽子の類の上のクーロン項は、最大を二重に閉じた点から動かす働き手として使える。最大の値は、すでに引用した入力のままである。

離散理論は、各固定コンパクト化において \(A_{\max}(N)<\infty\) を証明し、その値は評価しない~\cite{dst-dio}。見積もり \(A_*\sim 300\) は、その論文と本文のヒューリスティックであり、魔法数列の経験的な端と層の予算とから読まれている。第\ref{sec:layer}節の節は、それを厳しくしない。\(\ell\) が \(A\) に依らなければ、液滴半径の核表面は、すでに \(A=1\) においてすべての節の外側にある。代わりに \(\ell\) が \(A^{1/3}\) のように尺度をもてば、表面と節との比はすべて \(A\) に依らず、やはり打ち切りは現れない。この枠組みでの核種図表の有限性は、ローター像の有限性であり、\(\gamma_s\) の性質ではない。

\section{共有される入口と、核分裂}
\label{sec:fission}

双対励起は、近接を固定した双対角への付加であり、~\eqref{eq:ansatz} に補な読みである。
\begin{equation}
\label{eq:dual-ansatz}
\alpha=\frac{\ell}{r},
\qquad
\beta=\frac{\ell}{r}+\psi,
\qquad
\psi\ge 0.
\end{equation}
\(\varepsilon=\psi/2\) とし、\(\tau\) を~\eqref{eq:tau} のとおりとすると、
\begin{equation}
\label{eq:gamma-psi}
\gamma_\psi(x)
  =\cos\varepsilon\,\gamma_s(x)-\sin\varepsilon\,\tau(x)
\end{equation}
である。外側の井戸 \(0<x<x_1\) では、\(\gamma_s\) と \(\tau\) はともに正である。\(\psi\in(0,\pi)\) では \(\varepsilon\in(0,\pi/2)\) かつ \(\cos\varepsilon>0\) であるから、\(\gamma_\psi\) は \(\gamma_s/\tau-\tan\varepsilon\) の符号をもつ。零点は
\begin{equation}
\label{eq:psi-star}
\psi_*(x)=2\arctan\frac{\gamma_s(x)}{\tau(x)}
\end{equation}
にある。商 \(\gamma_s/\tau\) は \((0,x_1)\) の上で狭義に減る。その導関数の分子は \(\gamma_s'\tau-\gamma_s\tau'\) である。\(\gamma_s'=-2\cosh x\sin x\) と \(\tau'=2\cosh x\cos x\) を代入すると、括弧 \(\sin x\,\tau+\gamma_s\cos x\) が残り、その括弧は \(\cosh x\) に等しいので、分子は \(-2\cosh^{2}x\) である。したがって \(x\) が \(0\) から \(x_1\) まで走るあいだ、\(\psi_*\) は \(\pi\) から \(0\) まで厳密に落ちる。励起 \(\psi\in(0,\pi)\) のそれぞれは、外側の井戸のちょうど一つの位相 \(x_\psi\) における臨界値である。

小さな位相では \(\tau(x)=2x+O(x^{3})\)、\(\gamma_s(x)=1+O(x^{2})\) であるから、
\[
\psi_*(x)=\pi-4x+O(x^{3})=\pi-\frac{2\ell}{r}+O\Bigl(\Bigl(\frac{\ell}{r}\Bigr)^{3}\Bigr)
\]
である。大きく離れた対を反転するのは、\(\pi\) に近い励起だけである。すでに節の近くにある対は、小さな励起で反転する。相対ラピディティは、大きな間隔でこれを真似できない。\(2\ell/\pi\) の外側では、因子~\eqref{eq:gamma-phi} があらゆる有限な \(\varphi\) で正のままである。

位相 \(x_\psi\) は、\(\psi\) を固定したときの間隔の吸引子である。\(x>x_\psi\) では \(\psi_*(x)<\psi\) であるから \(\gamma_\psi<0\) であり、力は外向きである。したがって間隔は増し、位相は \(x_\psi\) へ向かって落ちる。\(x<x_\psi\) では力は内向きであり、間隔は縮み、位相は \(x_\psi\) へ向かって上がる。\((\gamma_s/\tau)'\neq 0\) であるから \(\gamma_\psi\) の零点は単純であり、力はそこに単純極をもつ。\(\psi\downarrow 0\) とすると吸引子は静的な節に近づく。\(\psi\uparrow\pi\) とすると \(x_\psi\downarrow 0\) であるから、吸引子は無限遠へ退く。先頭の不足から、\(\pi\) に近い励起における吸引子の半径は \(2\ell/(\pi-\psi)\) である。

ラピディティの閾値と双対の閾値は、一つの反応の二つの名前ではない。ラピディティの閾値 \(\varphi_*\) は入口であり、共有される。捕獲、中性子誘起核分裂、核融合は、それぞれ、なお \(\varphi_*\) の内向きの側にいる相手を要する。その閾値より上では核のチャンネルが斥け、複合系が存在するより前に相手は引き返される。核融合とは、この入口が完了したものである。合体した系は内向きのままであり、質量数は増えている。二つの荷電核に対しては第\ref{sec:proton}節のクーロン項が残り、二つの項の比較は第\ref{sec:cold}節の入口である。中性子はすでに電荷を担わない。それが遅いことが、\(\varphi\) を \(\varphi_*\) の下に保つ。

核分裂はその入口ではない。すでに形成された核の双対励起であり、高い相対ラピディティでも、結合そのものでもない。本文はそれを、\(\gamma_s\) を正から負へ駆動し、破片を斥ける励起された双対モードとして述べている~\cite{dst-main}。式~\eqref{eq:psi-star} がその遷移の位置を定める。\(x_\psi\) の内側では力は外向きであり、それが破片の斥力である。外側では外側の井戸がなお引力であるから、固定された励起は間隔を無限遠へ送るのではなく、\(r=\ell/(2x_\psi)\) に集める。逃げ去る分裂には、励起が \(\pi\) に近づき、集める半径が発散すること、あるいは分離のあいだに励起が変わることが要る。\(Q\) 値も、アクチノイド核分裂の非対称な質量分裂も、計算しない。中性子誘起核分裂は、両方の閾値をその順序で用いる。中性子は \(\varphi\) が \(\varphi_*\) より下で到着するか、あるいは決して加わらない。その結合が置く励起 \(\psi\) が、\(\psi_*\) を横切りうるものである。自発核分裂には中性子の入口がない。それは第二の閾値だけである。捕獲、散乱、核融合、核分裂のあいだの分岐は、\(\gamma\) の符号だけでは決まらない。

したがって二つの閾値は二つの異なる問いに答え、遅い中性子が両方に座るのは、入り方としてだけである。ラピディティの閾値は、すでに \(2\ell/\pi\) の内側の間隔にある動く核子が、チャンネルを外向きと見いだすときを言う。双対の閾値は、静止した対を、どの間隔でも、内部励起が反転するときを言う。結合は第一を用いる。分裂は第二を用いる。

\section{ラザフォード散乱}
\label{sec:rutherford}

金に当たるアルファ粒子は、数十フェムトメートルの間隔で~\eqref{eq:two-channel} のクーロン項に出会う。ラザフォードの法則 \(d\sigma/d\Omega=(d/4)^{2}/\sin^{4}(\theta/2)\)、\(d=2Zke^{2}/E\) は、正面衝突の距離 \(d\) まで符号のないクーロン斥力が続くときの法則である~\cite{rutherford1911,geiger1913}。入力 \(Z=79\)、\(E=7.69\,\mathrm{MeV}\)、\(ke^{2}=1.439964\,\mathrm{MeV\,fm}\) のとき、実験室系の距離 \(d\) は \(29\) と \(30\,\mathrm{fm}\) のあいだに厳密にある。

第~\ref{rule:length} 節の長さの規則により、アルファ粒子と金の原子核は、陽子が陽子のカオスにある核子の集まりとして相互作用し、そのあいだで読む長さは、~\eqref{eq:two-channel} のクーロン項についても \(\lambda_\pi\) である。これは \(\ell_E=\ell=\lambda_\pi\) の場合である。正面衝突の距離は \(20\lambda_\pi\) を超え、第一の節の外にある。\(29\,\mathrm{fm}\) 以上のすべての間隔では、符号つきの因子が一を上回る分は一万分の六未満であり、ずれは \((\lambda_\pi/2r)^{2}\) のように増える。

重心系の運動エネルギーは、標的の質量数を \(A\) として、実験室のエネルギーの \(A/(A+4)\) 倍である。したがって正面衝突の距離は \((A+4)/A\) 倍になり、そこでのずれは \(5.5\times 10^{-4}\) である。符号つきの力の古典的な偏向は、断面積がラザフォードの値を上回る分として次を与える。\(7.69\,\mathrm{MeV}\) では \(90^\circ\) で \(3.1\times 10^{-4}\)、\(170^\circ\) で \(5.8\times 10^{-4}\)。\(5.3\,\mathrm{MeV}\) では \(1.5\times 10^{-4}\) と \(2.8\times 10^{-4}\)。後方散乱分析のエネルギーである \(2\,\mathrm{MeV}\) では \(4\times 10^{-5}\) 未満にとどまる。これらは、大角度のデータがラザフォードの法則と比べられる一パーセントの精度よりはるかに小さい。同じ力のポテンシャルを、娘核とアルファ粒子の接触半径 \(1.2(A_d^{1/3}+4^{1/3})\,\mathrm{fm}\) から転回点まで積分すると、\({}^{212}\mathrm{Po}\) のガモフ指数は \(0.037\)、\({}^{238}\mathrm{U}\) のそれは \(0.033\) だけ上がる。半減期は \(3.8\) および \(3.4\) パーセント延びるが、これは前形成因子の幅のはるか内側である。

重い原子核によるアルファ粒子の異常散乱は、二つの表面が近づくところ、金では \(1.2\,(A_1^{1/3}+A_2^{1/3})\,\mathrm{fm}\approx 9\,\mathrm{fm}\) 付近で始まる。\(10\,\mathrm{fm}\) で因子が一からずれるのは \(5\times 10^{-3}\) である。この読みでは、すきまを挟んだ核子の対が \(\lambda_\pi\) の数倍以内に入り、それぞれの間隔で節に近づくときに、原子核の大きさのずれが始まる。そこはデータがずれる場所である。式~\eqref{eq:two-channel} は二体の力であり、それらの対を足し合わせてはいない。

核の前因子 \(\kappa\) は対の入力であり、二つの原子核の核子にわたってそれがどう足し合わされるかを本論文は定めない。\(A_1A_2\) 個の対にわたってコヒーレントに足されるなら、遠方の前因子は \(Z_1Z_2ke^{2}-A_1A_2\kappa\) となり、断面積は \((1-\eta)^{2}\) 倍される。ここで \(\eta=A_1A_2\kappa/(Z_1Z_2ke^{2})\) である。金とアルファ粒子では \(A_1A_2/(Z_1Z_2)\) は \(4.98\) と \(4.99\) のあいだに厳密にある。遠方での割合が \(1/200\) 以下、つまり断面積の一パーセントの変化以下であるなら、\(\kappa<ke^{2}/997\) である。ニュートンの前因子 \(Gm_p^{2}\) では \(\eta\approx 4\times 10^{-36}\) である。ここで長さは質量数とともに変わらない。\(\lambda_\pi\) は核子の対の長さである。

電荷とともに断面積が増えるのは、クーロン項の増え方である。ラザフォードの法則は \(Z^{2}\) に比例する。チャドウィックは、銅、銀、白金でのアルファ粒子の散乱から、原子番号 \(29\)、\(47\)、\(78\) に対して核電荷 \(29.3\)、\(46.3\)、\(77.4\) を、約一パーセントの誤差で得た~\cite{chadwick1920}。コヒーレントな和のもとでは、有効電荷は \(Z(1-\eta)\) であり、アルファ粒子では \(\eta=2(A/Z)\,\kappa/(ke^{2})\) である。係数 \(2A/Z\) は、銅で \(4.38\) と \(4.39\)、銀で \(4.59\) と \(4.60\)、白金で \(5.00\) と \(5.01\) のあいだにある。三つの電荷をそれぞれ一パーセントの誤差で最小二乗法により合わせると、\(\kappa/(ke^{2})=(0.94\pm 1.24)\times 10^{-3}\) となり、零と矛盾しない。重力の項は質量数とともに増え、測られた増え方は電荷に従う。\(\lambda_\pi\) での符号つきのクーロン因子は、電荷とともに減る小さな超過を加える。軽い原子核ほど近くまで接近されるからである。同じ重心系の \(7.69\,\mathrm{MeV}\)、\(150^\circ\) での超過は、アルミニウムで \(1.6\times 10^{-2}\)、銅で \(3.8\times 10^{-3}\)、銀で \(1.5\times 10^{-3}\)、金で \(5.5\times 10^{-4}\) である。アルミニウムではその経路は \(5.7\,\mathrm{fm}\) に達し、接触半径 \(1.2\,(A_1^{1/3}+A_2^{1/3})\,\mathrm{fm}\approx 5.5\,\mathrm{fm}\) にかかる。そこでは二体の式だけでは記述が尽きない。

金では、この散乱の間隔はその上限の内側に留まる。残る入力は、核子にわたる \(\kappa\) の和である。

\section{中性子の散乱}
\label{sec:neutron}

中性子は重力チャンネルだけを担う。核子に及ぼす力は、長さ \(\lambda_\pi\) での~\eqref{eq:force} に第~\ref{sec:rapidity} 節の相対ラピディティを加えたものであり、その斥力は~\eqref{eq:gamma-phi} を通して速さとともに変わる。斥力の始まりは、遅い対では静的な節 \(r_1=0.7540\,\mathrm{fm}\) にあり、ラピディティが増すにつれて外へ動く。\(0.0253\,\mathrm{eV}\) の熱中性子ではラピディティは \(7.34\times 10^{-6}\) で、始まりは千万分の \(9\) だけ動く。\(1\,\mathrm{MeV}\) ではラピディティは \(0.0461\) で、始まりは \(0.7582\,\mathrm{fm}\) にある。\(14\,\mathrm{MeV}\) ではラピディティは \(0.172\) で、始まりは節の二パーセント外の \(0.7695\,\mathrm{fm}\) にある。どの速さでも、それは壁 \(2\lambda_\pi/\pi=0.9001\,\mathrm{fm}\) の内側にとどまる。したがって各核子の対は、中性子のエネルギーとともに厚くなる斥力の皮をもち、そのエネルギー依存性は重力チャンネルそのものの性質である。

測られた散乱は、どの原子核でも同じというわけではない。熱中性子の束縛干渉性散乱長は、アルミニウムで \(3.449\,\mathrm{fm}\)、銅で \(7.718\)、銀で \(5.922\)、金で \(7.63\)、鉛で \(9.405\) であり、水素 \(-3.739\)、チタン \(-3.438\)、バナジウム \(-0.3824\,\mathrm{fm}\) では負である~\cite{sears1992}。全束縛散乱断面積は、アルミニウムの \(1.503\,\mathrm{b}\) から鉛の \(11.118\,\mathrm{b}\)、ニッケルの \(18.5\,\mathrm{b}\) まで、\(A\) について単調でない順に並ぶ。\(14\,\mathrm{MeV}\) での評価済み全断面積は、\({}^{27}\mathrm{Al}\) で \(1.744\,\mathrm{b}\)、\({}^{63}\mathrm{Cu}\) で \(2.909\)、\({}^{65}\mathrm{Cu}\) で \(2.959\)、\({}^{107}\mathrm{Ag}\) で \(4.258\)、\({}^{109}\mathrm{Ag}\) で \(4.396\)、\({}^{197}\mathrm{Au}\) で \(5.295\)、\({}^{208}\mathrm{Pb}\) で \(5.385\,\mathrm{b}\) である~\cite{endf8}。幾何学的な形 \(2\pi(r_0A^{1/3}+\hbar/p)^{2}\) は、\(r_0\) を \(1.35\) から \(1.50\,\mathrm{fm}\) のあいだにとると、これらの \(14\,\mathrm{MeV}\) の断面積を再現し、この区間は \(\lambda_\pi\) を含む。\(14\,\mathrm{MeV}\) での \(\hbar/p\) は \(1.2\,\mathrm{fm}\) である。本稿はその一致を導出しない。

これらの領域のあいだで、換算波長は熱エネルギーの \(2.9\times 10^{4}\,\mathrm{fm}\) からその \(1.2\,\mathrm{fm}\) まで下がり、データのエネルギー依存性はそれと各原子核の共鳴とに従う。金の評価済み全断面積は \(1\,\mathrm{MeV}\) の \(5.755\,\mathrm{b}\) から \(14\,\mathrm{MeV}\) の \(5.295\,\mathrm{b}\) へ下がり、そのあいだに各対の皮は一・五パーセント広がる。皮は核子の対の尺度での補正であり、測られた変化の大きな部分は原子核の大きさと換算波長が担う。遠方は別の補正であり、それは定まる。位相が消えるところでは、対の干渉因子は \(\cosh(\varphi/2)\) に等しく、力は強さ \(\kappa/\cosh(\varphi/2)\) の逆二乗則である。中性子の静止エネルギー \(939.565\,\mathrm{MeV}\) を用いると、半角の変形はその強さを \(\kappa/\sqrt{1+T/(2m_n c^{2})}\) と書く。\(1\,\mathrm{MeV}\) ではこの平方根は \(1.000266\) と \(1.000267\) のあいだに厳密にある。\(14\,\mathrm{MeV}\) では \(1.003718\) と \(1.003719\) のあいだに厳密にある。前者から後者へ、係数が減る割合は \(4/1000\) より厳密に小さい。その減少は皮の広がりより小さく、金の全断面積の変化を担わない。

位相 \(x\in(0,\pi/4)\) のあらゆる点、あらゆる有限なラピディティで、力は内向きのままである。この逆二乗則を上回る割合は、\(e^{-x}/(\cos x-\sin x)-1\) より小さい。位相が高々 \(7/100\) のとき、その割合は \(6/1000\) より厳密に小さい。

中性子がすべての核子のラピディティの壁の外にいるあいだ、対の力の和は中性子を引き返させない。核子はある球の中にある。中性子が、\(2\lambda_\pi/\pi\) だけ広げた同心球の外にいるなら、どの対も自分の壁の外にある。各対の力はその核子へ向かい、正味の力の動径成分は負である。球の内部の配置も、有限な速さも、壁に達する前にその成分を反転させない。引き込みのコヒーレントな部分は、強さ \(A\kappa/\cosh(\varphi/2)\) のラザフォード場である。目立つ大きさの \(\kappa\) は、\(A^{2}\) 倍された、その形の前方の山を加える。ニュートンの前因子 \(Gm_p^{2}\) では、その山は無視できる。速さに依存する皮はこの場にはない。皮が始まるのは、ある対が自分の壁に入るときだけである。原子核ごとの違い、負の散乱長を含めて、それらの遭遇に属する。その振幅は、ここでは導出しない。

\section{結論}
\label{sec:conclusions}

斥力反転層は、等スケール因子 \(\gamma_s\) が負である最初の区間である。その外側の壁は節 \(x_1=0.93755\) であり、\(r_1/\ell=0.53330\) にあり、内側の壁は \(r_2/\ell=0.12731\) である。外側の壁は、両側から間隔を引き寄せる。等スケール軌跡の上では、間隔はその壁へ向けて駆動され、層を貫かない。自由な落下は、壁の厳密に外側でラピディティの閾値に出会い、層には入らない。層の内部で、\(\gamma_s\) の極小における外向きの力の平坦部の値は \(3.4057\,\kappa/\ell^{2}\) である。前因子 \(\kappa\) は核の入力である。組合せ \(G m_p^2\) は \(k e^2\) より \(1.236\times 10^{36}\) だけ小さく、その入力ではない。長さは荷電パイオンのコンプトン波長 \(\lambda_\pi\) である。実験室の符号は、それを \(56/75\,\mathrm{fm}\) と \(19/10\,\mathrm{fm}\) のあいだに厳密に置き、その第一の節は芯の尺度 \(\tfrac12\,\mathrm{fm}\) と引力の半径 \(1\,\mathrm{fm}\) のあいだに厳密にある。

通常の近接ラピディティに加えられた相対ラピディティ \(\varphi\) は、干渉因子を下げる。位相 \(x\in(\pi/4,x_1)\) での力は、\(\varphi<\varphi_*(x)\) で内向き、\(\varphi>\varphi_*(x)\) で外向きであり、\(\varphi_*\) は閉じた対数~\eqref{eq:phi-star} である。有限なラピディティは \(2\ell/\pi\) の外側で力を反転せず、壁の外向きの変位は比 \(1.1937\) で抑えられる。遅い中性子の斥力の表皮は、分数的な厚さ \(0.12295\,\varphi\) をもつ。\(1\,\mathrm{MeV}\) では壁は千分の六だけ動き、\(100\,\mathrm{MeV}\) では五パーセント動く。自身の閾値の内向きの側にいる中性子は引き込まれる。外向きの側にいる中性子は、極から押し出される。内向きのラピディティが増え続ければ、引き込みは閾値の上で終わり、静的な節の厳密に外側である。その二者択一が、捕獲と、中性子誘起核分裂と、核融合とに共有される入口である。ラピディティを捨てることと、\(1/v\) の法則とは、導出されない。

二つの陽子は、符号つきのクーロン力を加える。\(\gamma_s\) が零を下向きに横切る壁は質量の対を引き寄せ、同符号の電荷の対を斥け、どちらのチャンネルでも隣り合う壁は逆の安定性をもつ。等しくない長さは壁を互い違いにするので、核の配置は符号胞である。一つの瞬間に内向きの対と外向きの対とがあり、各壁の上の力は非有界である。その胞がカオス的な陽子力である。等しい長さは、そのような互い違いを生まず、単一の前因子 \(k e^{2}-\kappa\) だけを生む。

閉包 \(2\)、\(8\)、\(20\) と、\(28\)、\(50\)、\(82\)、\(126\) とは、累積の振動子数 \(C(N)\) と、ずらした数 \(C(N)+2(N+2)\) である。それらは各ローター類の席をラベルし、節をラベルしない。二重魔法核種は両方の類を閉じる。安定性、谷、および \(A=62\) 付近の \(\mathrm{BE}/A\) の最大は、閉包の帰結ではなく、\(A_{\max}\) は離散理論の、有限だが未評価の限界のままである。

双対励起 \(\psi\) は外側の井戸を一つの位相で反転し、その位相は間隔を引き寄せる。臨界角は、無限遠の \(\pi\) から第一の節の \(0\) まで落ちる。核分裂は、すでに形成された核の、この反転である。核融合はその反転ではない。共有される入口 \(\varphi<\varphi_*\) であり、合体した系がなお内向きであるあいだに完了する。遅い中性子が核分裂を誘起するのは、その入口を通過するときだけである。分裂が続くならば、それは \(\psi>\psi_*\) である。相対ラピディティは、\(2\ell/\pi\) の外側で分裂を生まない。

荷電の対に対して、同じ入口は極の遅い側である。\(\varphi_*\) より上では核の項がクーロン項に加わり、正味の斥力はクーロン斥力を上回る。\(\varphi_*\) より下では核の項がそれに対抗し、電磁の因子が正であるかぎり、閾値のすぐ下ではその対抗が、あらゆる有限な比 \(k e^{2}/\kappa\) を上回る。その側で最も強い内向きの力は、静止ではなく極のすぐ下にある。より低いラピディティは、極を節の近くへ動かし、静止エネルギーの高々 \(10^{-10}\) の運動エネルギーは、半径を \(5\times 10^{-6}\) より少なく広げる。核子または重陽子の室温は、その類に入る。\(2\ell/\pi\) の外側では、大きなラピディティが核の引力を消し、クーロン斥力を残す。自由な落下はなお節より前に極に出会うので、低温の入口は、ラピディティが抑えられたままの経路である。その保持を一軸に沿って引き受ける相手は、正イオンの正八面体において、第一の電荷の節の外側にある戻す席である~\cite{dst-cold}。反応率は導出されない。強い力は、全体を通して、一つの干渉因子の符号である。第一の節の外側では内向き、反転層では外向き、古典的な落下が達しないより深い室でのみ再び内向き、そして \(\varphi_*\) を越えたラピディティか \(\psi_*\) を越えた双対励起かのどちらかのもとで外向きである。

すべての核子のラピディティの壁の外にいる中性子は、内へ引かれる。各対の遠方は、強さ \(\kappa/\cosh(\varphi/2)\) の逆二乗則であり、位相が高々 \(7/100\) のとき、それを上回る分は \(6/1000\) より小さい。コヒーレントな引き込みは、強さ \(A\kappa/\cosh(\varphi/2)\) のラザフォード場である。\(1\,\mathrm{MeV}\) から \(14\,\mathrm{MeV}\) へ、その係数が減る割合は \(4/1000\) より厳密に小さい。皮が中性子に出会うのは壁の内側だけであり、その遭遇の振幅は導出されない。

\begin{thebibliography}{9}

\bibitem{dst-main}
Anonymous,
``Gravity as Torsion between Dual Spacetime: A Biquaternionic Reformulation of General Relativity'' (2025).
\url{https://github.com/hypernumbernet/dual-spacetime-doc/blob/main/double-spacetime-theory.pdf}.

\bibitem{dst-shell}
Anonymous,
``Electron Orbits from the Signed Electrostatic Force: Covalent Bonds, Metallic Bonds, and Crystals,''
companion note.

\bibitem{dst-cold}
Anonymous,
``Cold Fusion in a Metal: A Restoring Seat in an Octahedral Cage,''
companion note.

\bibitem{dst-dio}
Anonymous,
``Discrete Biquaternionic Double Spacetime: Rigorous Proofs of Diophantine Conjectures via Torsional Mismatch and PGA Embedding,''
companion note.

\bibitem{codata2018}
E.~Tiesinga, P.~J.~Mohr, D.~B.~Newell, and B.~N.~Taylor,
``CODATA recommended values of the fundamental physical constants: 2018,''
\textit{Rev.\ Mod.\ Phys.} \textbf{93}, 025010 (2021).

\bibitem{wiringa1995}
R.~B.~Wiringa, V.~G.~J.~Stoks, and R.~Schiavilla,
``Accurate nucleon-nucleon potential with charge-independence breaking,''
\textit{Phys.\ Rev.\ C} \textbf{51}, 38 (1995).

\bibitem{angeli2013}
I.~Angeli and K.~P.~Marinova,
``Table of experimental nuclear ground state charge radii: An update,''
\textit{At.\ Data Nucl.\ Data Tables} \textbf{99}, 69 (2013).

\bibitem{rutherford1911}
E.~Rutherford,
``The scattering of \(\alpha\) and \(\beta\) particles by matter and the structure of the atom,''
\textit{Phil.\ Mag.} \textbf{21}, 669 (1911).

\bibitem{geiger1913}
H.~Geiger and E.~Marsden,
``The laws of deflexion of \(\alpha\) particles through large angles,''
\textit{Phil.\ Mag.} \textbf{25}, 604 (1913).

\bibitem{chadwick1920}
J.~Chadwick,
``The charge on the atomic nucleus and the law of force,''
\textit{Phil.\ Mag.} \textbf{40}, 734 (1920).

\bibitem{sears1992}
V.~F.~Sears,
``Neutron scattering lengths and cross sections,''
\textit{Neutron News} \textbf{3}, No.~3, 29 (1992).

\bibitem{endf8}
D.~A.~Brown \textit{et al.},
``ENDF/B-VIII.0: The 8th major release of the nuclear reaction data library,''
\textit{Nucl.\ Data Sheets} \textbf{148}, 1 (2018).

\end{thebibliography}

\end{document}
